\chapter{Experimental Results}
\label{chap:results}

Every ordinary panel in this chapter begins at a true MI difference of zero.
That point measures false-positive rejection under \(H_0\). Points to its
right measure detection under \(H_1\). The horizontal coordinate is a
\emph{constructed population} difference in nats. The vertical coordinate is
the rejection rate over 20,000 random table pairs. The blue solid circles are Normal Wald and the
magenta dashed squares are Expanded Welch. The dotted horizontal line marks
the nominal 0.05 rejection rate. Each plotted rate divides the number of valid rejections
by all 20,000 replicates. A hollow marker indicates a valid-result rate below
0.90. The figure source manifest gives the exact saved
configuration identifiers for every panel.

\section{Main landscape: shape, margins and sample size}

The first question is broad: do the methods behave differently as the table
grows, margins become uneven and samples become larger?
Figures~\ref{fig:main-2x2-uniform}--\ref{fig:main-8x8-different_skew} provide a
systematic view, not a selection of favourable examples. Each of the
four shapes is shown under all three margin profiles. Within every figure,
the columns are the same three equal sample sizes. The full vertical scale
preserves comparability of detection; the lower row enlarges the region in
which calibration differences near 0.05 can be seen. The full nine-size
ladder, including \(n=2,5,20,50,250,500\), remains in the companion atlas.

\newcommand{\resultinterpretation}[1]{%
\par\smallskip\noindent\begin{minipage}{\textwidth}
\small\textbf{Interpretation.} #1
\end{minipage}%
}

\newcommand{\mainlandscapefigure}[6]{%
\begin{figure}[p]
\centering
\makebox[\textwidth][c]{\includegraphics[width=1.18\textwidth]{main_#1_#3.pdf}}
\caption[Main landscape: #2 tables, #5 margins]{#2 tables, #4. Columns: \(n_P=n_Q=10,100,1000\), respectively.
For each population the first two rows and columns receive the signed changes
in Equation~\eqref{eq:first-change-block}; \(I(P)=0.02\) and
\(I(Q)=0.02+|\Delta|\). The horizontal grid is
\(|\Delta|=0,0.001,0.002,0.005,0.01,0.02\) nats. The upper panels have a
shared 0--1 rejection axis and the lower panels show the same data at 0--0.15.
At zero the rate is false positives; elsewhere it is detection. Each point
uses 20,000 paired samples. Hollow markers denote validity below 90\%.}
\label{fig:main-#1-#3}
\resultinterpretation{#6}
\end{figure}%
}

\mainlandscapefigure{2x2}{\(2\times2\)}{uniform}{both \(P\) and \(Q\) have uniform row and column margins \((1/2,1/2)\)}{uniform}{At zero difference, both methods reject below 0.05 at \(n=100\) and 1,000; Expanded Welch is lower still. At \(n=10\), Wald is slightly above 0.05 and Expanded Welch below it. Validity is high in all three panels. Rejection grows most clearly with MI difference at \(n=1{,}000\), where the two curves are close. At the smaller sizes, low rejection should not be mistaken for better detection.}
\mainlandscapefigure{2x2}{\(2\times2\)}{same_skew}{both \(P\) and \(Q\) have row and column margins \(m_2(0.8)=(0.8,0.2)\)}{same skew}{The null points are below 0.05, especially for Expanded Welch. At \(n=10\), many Expanded Welch calculations are invalid (hollow markers), so its low unconditional rate has more than one cause. Detection rises with the constructed difference mainly at \(n=1{,}000\); there Wald remains slightly higher. This skewed pair gives no general calibration advantage to the more conservative reference.}
\mainlandscapefigure{2x2}{\(2\times2\)}{different_skew}{\(P\) has row and column margins \(m_2(0.7)=(0.7,0.3)\), while \(Q\) has \(m_2(0.8)=(0.8,0.2)\)}{different skew}{At zero difference, the two tables have equal MI but different margins. Both methods are mostly below the 0.05 target, and Expanded Welch is lower. At \(n=10\), its hollow markers also indicate substantial invalidity. The alternative curves rise noticeably at \(n=1{,}000\), with Wald just above Expanded Welch. The same conservative pattern therefore persists when the equal-MI null does not require identical populations.}

\mainlandscapefigure{3x3}{\(3\times3\)}{uniform}{both \(P\) and \(Q\) have uniform row and column margins \((1/3,1/3,1/3)\)}{uniform}{At \(n=10\), Wald's zero-difference rate is above 0.05 and Expanded Welch brings it below the line; both methods are almost always valid. At \(n=100\), however, both start well below 0.05 and the correction lowers rejection further. Detection grows strongly at \(n=1{,}000\), where Wald remains slightly higher. Whether the adjustment helps calibration changes with sample size, not just table shape.}
\mainlandscapefigure{3x3}{\(3\times3\)}{same_skew}{both \(P\) and \(Q\) have row and column margins \(m_3(0.8)=(0.8,0.1,0.1)\)}{same skew}{The \(n=10\) null rate falls from slightly above 0.05 for Wald to below it for Expanded Welch, but Expanded Welch is valid in only about two thirds of those samples. At \(n=100\) and 1,000, both null rates are below 0.05. Rejection rises with the MI difference most visibly at \(n=1{,}000\); Expanded Welch remains lower. The small-sample null reduction is not solely a reference-distribution effect.}
\mainlandscapefigure{3x3}{\(3\times3\)}{different_skew}{\(P\) has row and column margins \(m_3(0.7)=(0.7,0.15,0.15)\), while \(Q\) has \(m_3(0.8)=(0.8,0.1,0.1)\)}{different skew}{With different margins, Wald rejects too often at the \(n=10\) null; Expanded Welch rejects less, although its hollow markers show reduced validity. At \(n=100\), both are conservative. At \(n=1{,}000\), detection increases with the true difference and the methods draw closer. This comparison shows a small-sample benefit against liberal rejection but not a consistent improvement across sample sizes.}

\mainlandscapefigure{5x5}{\(5\times5\)}{uniform}{both \(P\) and \(Q\) have uniform row and column margins \((1/5,\ldots,1/5)\)}{uniform}{Both methods are valid in the displayed panels. At zero difference, Wald is liberal at \(n=10\), while Expanded Welch reduces but does not eliminate the excess. At \(n=100\) and 1,000, both null rates are below 0.05. Little detection is gained over the displayed difference range at \(n=100\), whereas rejection grows clearly at \(n=1{,}000\). The correction lowers both false positives and detection.}
\mainlandscapefigure{5x5}{\(5\times5\)}{same_skew}{both \(P\) and \(Q\) have row and column margins \(m_5(0.8)=(0.8,0.05,0.05,0.05,0.05)\)}{same skew}{At \(n=10\), Wald starts above the 0.05 line and Expanded Welch below it, but the latter is invalid for roughly one third of samples. At \(n=100\) and 1,000, both begin below the line. Detection increases substantially only in the largest-sample panel, with Wald higher. The small-sample appearance of better calibration must be read together with validity.}
\mainlandscapefigure{5x5}{\(5\times5\)}{different_skew}{\(P\) has row and column margins \(m_5(0.7)=(0.7,0.075,0.075,0.075,0.075)\), while \(Q\) has \(m_5(0.8)=(0.8,0.05,0.05,0.05,0.05)\)}{different skew}{The \(n=10\) null is liberal for both methods, though less so for Expanded Welch; its hollow markers also show reduced validity. At \(n=100\), rejection stays low and even falls across this short alternative range, so these curves should not be described as increasing power. At \(n=1{,}000\), detection rises and Wald remains slightly higher. The benefit of reducing excess null rejection is regime-specific.}

\mainlandscapefigure{8x8}{\(8\times8\)}{uniform}{both \(P\) and \(Q\) have uniform row and column margins \((1/8,\ldots,1/8)\)}{uniform}{At \(n=10\) and 100, the zero-difference rates exceed 0.05 for both methods; Expanded Welch reduces the excess but does not always reach the target. These plotted calculations are valid, so the difference is not caused by missing outputs. At \(n=100\), rejection changes little over the displayed alternatives; at \(n=1{,}000\), it grows strongly. Even with uniform margins, table size and sample size alter the trade-off.}
\mainlandscapefigure{8x8}{\(8\times8\)}{same_skew}{both \(P\) and \(Q\) have row and column margins \(m_8(0.8)=(0.8,0.2/7,\ldots,0.2/7)\)}{same skew}{At \(n=10\), Wald is above 0.05 at the null and Expanded Welch below it, but many Expanded Welch outputs are invalid. Both null rates are below 0.05 at \(n=100\) and 1,000. Across the alternative points, appreciable growth appears chiefly at \(n=1{,}000\), with Wald higher. The adjustment does not turn the smaller-sample curves into strong detection.}
\mainlandscapefigure{8x8}{\(8\times8\)}{different_skew}{\(P\) has row and column margins \(m_8(0.7)=(0.7,0.3/7,\ldots,0.3/7)\), while \(Q\) has \(m_8(0.8)=(0.8,0.2/7,\ldots,0.2/7)\)}{different skew}{At zero difference, both methods are too liberal at \(n=10\) and 100, although Expanded Welch rejects less. Its \(n=10\) markers are hollow, so validity contributes to that reduction. At \(n=100\), rejection falls rather than grows over the displayed alternative range; at \(n=1{,}000\), detection rises but stays lower for Expanded Welch. This is the clearest reminder not to infer power from the null rate alone.}

\clearpage

The landscape does not yield a uniformly preferable method. At \(n=5\),
Expanded Welch has a zero null rejection rate in all
\(\MainFiveNullCount\) main regimes. This includes the
uniform \(3\times3\), \(5\times5\) and \(8\times8\) regimes, where more than
90\% of calculations are valid, so the zero rates cannot be attributed only
to invalid outputs. Among the \(\MainTenTwentyNullCount\) main null regimes with \(n=10\) or 20,
Expanded Welch is numerically closer to 0.05 in
\(\MainTenTwentyExpandedCloserCount\) and Wald in
\(\MainTenTwentyWaldCloserCount\). The comparison reverses
in the \(\MainLateNullCount\) main null settings at \(n=500\) or 1,000: Wald
is closer to 0.05 in all \(\MainLateWaldCloserCount\). At \(n=1{,}000\), Wald
null rates across the main settings lie between
\(\MainThousandWaldMinRate\) and \(\MainThousandWaldMaxRate\); Expanded
Welch lies between \(\MainThousandExpandedMinRate\) and
\(\MainThousandExpandedMaxRate\). Neither is uniformly at the nominal
level.

Expanded Welch uses a heavier-tailed reference and can only lower rejection
relative to Wald on the same table pairs. This can bring a liberal null rate
toward 0.05, but moves a conservative rate farther below it. Closeness to
0.05 must therefore be judged regime by regime.

Table~\ref{tab:null-paired-examples} puts two fully valid null points side by
side to show the two directions of this effect.

\begin{table}[htbp]
\centering
\small
\caption{Two fully valid main-grid null regimes. Brackets beside each rate
are pointwise Wilson 95\% intervals; the final bracket is the approximate
paired 95\% interval for Wald minus Expanded Welch on the same samples.
Source configurations are \texttt{config\_313a5e5c05280f9e} and
\texttt{config\_700945d218841f2e}, respectively.}
\label{tab:null-paired-examples}
\begin{tabularx}{\textwidth}{@{}l>{\centering\arraybackslash}X>{\centering\arraybackslash}X>{\centering\arraybackslash}X@{}}
\toprule
Regime & Wald & Expanded & Wald minus Expanded\\
 & Rate [95\% interval] & Rate [95\% interval] & Rate [95\% interval]\\
\midrule
\(8\times8\), uniform, \(n=20\) & 0.0687 [0.0653, 0.0723] & 0.0570 [0.0538, 0.0602] & 0.0118 [0.0103, 0.0132]\\
\(2\times2\), uniform, \(n=1000\) & 0.0445 [0.0417, 0.0474] & 0.0387 [0.0361, 0.0414] & 0.0059 [0.0048, 0.0069]\\
\bottomrule
\end{tabularx}
\end{table}

\FloatBarrier

In the uniform \(8\times8\), \(n=20\) regime, Expanded Welch moves Wald's
above-target rate nearer to 0.05. In the uniform \(2\times2\), \(n=1{,}000\)
regime, it moves an already-below-target rate farther from 0.05. Both methods
are valid on all replicates in these examples, so these changes reflect their
different reference thresholds, not missing outputs.

The \(n=5\) records in the companion atlas also show an important limit to
interpreting small rejection rates. In the \(8\times8\) different-skew
alternative with
\(n_P=n_Q=5\), \(I(P)=0.02\) and \(I(Q)=0.021\), Wald rejects
\(\EightSparseFiveAlternativeWaldRate\) and Expanded Welch rejects
\(\EightSparseFiveAlternativeExpandedRate\), but their valid-result rates
are only \(\EightSparseFiveAlternativeWaldValid\) and
\(\EightSparseFiveAlternativeExpandedValid\). These are probabilities of a
\emph{valid rejection}, not a clean comparison of power at matched achieved
false-positive rates. The corresponding null and validity records must be
examined before calling the gap a power advantage.

\section{Baseline MI and the location of changed cells}

The main grid fixes \(I(P)=0.02\). Does starting closer to independence
change the conclusion? Figure~\ref{fig:baseline-focus} holds shape, margins,
sample size and changed cells fixed while varying only the common starting
MI\@.

\begin{figure}[htbp]
\centering
\includegraphics[width=0.96\textwidth]{baseline_3x3_different_skew.pdf}
\caption[Baseline-MI comparison, \(3\times3\) tables]{Baseline-MI comparison for \(3\times3\) tables. Columns have
\(I(P)=0.0001,0.001,0.02\) nats, respectively; \(I(Q)=I(P)+|\Delta|\).
\(P\) has row and column margins \((0.7,0.15,0.15)\) and \(Q\) has
\((0.8,0.1,0.1)\). Both use the first-cell change block,
\(n_P=n_Q=100\), and \(|\Delta|=0,0.001,0.002,0.005,0.01,0.02\) nats.
Top: 0--1 rejection scale; bottom: 0--0.15 zoom. Each point has 20,000
replicates; hollow markers flag validity below 90\%.}
\label{fig:baseline-focus}
\resultinterpretation{For \(I(P)=0.0001,0.001,0.02\), the zero-difference Wald rates are
\(\BaselineTinyWaldRate\), \(\BaselineSmallWaldRate\) and
\(\BaselineRegularWaldRate\); Expanded Welch gives
\(\BaselineTinyExpandedRate\), \(\BaselineSmallExpandedRate\) and
\(\BaselineRegularExpandedRate\). All these null calculations are valid.
Rejection rises only modestly across the displayed alternatives, and both
methods remain mostly below 0.05. Starting closer to independence does not
make the larger Student cutoff helpful here: it lowers an already-conservative
null rate and further reduces detection.}
\end{figure}

Figure~\ref{fig:patterns-focus} asks whether the association being placed in
an early, rare or distributed set of cells changes performance when the
margins and target MI are fixed. The full \(3\times3\), \(5\times5\) and
\(8\times8\) matched sets are in the atlas.

\begin{figure}[htbp]
\centering
\includegraphics[width=0.96\textwidth]{patterns_5x5_different_skew.pdf}
\caption[Changed-cell comparison, \(5\times5\) tables]{Changed-cell comparison for \(5\times5\) tables. Columns place
association in the first \(2\times2\) block, last (rare) \(2\times2\)
block, or a centred score outer product spread across all cells. For every
column, \(P\)'s row and column margins are
\((0.7,0.075,0.075,0.075,0.075)\), and \(Q\)'s are
\((0.8,0.05,0.05,0.05,0.05)\). Here \(I(P)=0.0001\),
\(I(Q)=I(P)+|\Delta|\), \(n_P=n_Q=100\), and
\(|\Delta|=0,0.0001,0.00025,0.0005,0.001,0.002\) nats.
Top: 0--1; bottom: 0--0.15. Each point has 20,000 replicates.}
\label{fig:patterns-focus}
\resultinterpretation{For first, rare and spread patterns, respectively, the
zero-difference Wald rates are \(\PatternFirstWaldRate\),
\(\PatternRareWaldRate\) and \(\PatternSpreadWaldRate\); Expanded Welch gives
\(\PatternFirstExpandedRate\), \(\PatternRareExpandedRate\) and
\(\PatternSpreadExpandedRate\). All displayed calculations are valid. The
alternative curves remain nearly flat over this short MI range, so neither
method detects these differences well. Changing the location of association
does not remove the conservative gap here, although this one plot does not
make all patterns equivalent.}
\end{figure}

\section{Unequal samples and the direction of the difference}

Figures~\ref{fig:imbalance-q} and~\ref{fig:imbalance-p} hold the unordered
sample sizes at 50 and 250 but reverse which population has the larger
sample. They also show the two directions of the MI difference separately.
The distinction matters because the finite bias corrections are
\((r-1)(c-1)/(2n_P)\) and \((r-1)(c-1)/(2n_Q)\), and the two populations
have different margins.

\begin{figure}[htbp]
\centering
\includegraphics[width=0.96\textwidth]{imbalance_3x3_q_higher.pdf}
\caption[Unequal samples, \(Q\) higher]{Unequal samples, \(Q\) higher. Columns have
\((n_P,n_Q)=(50,250)\) and \((250,50)\). For both columns, \(P\) has
\(3\times3\) row and column margins \((0.7,0.15,0.15)\) and \(Q\) has
\((0.8,0.1,0.1)\); the first-cell change block is used.
\(I(P)=0.02\) and \(I(Q)=0.02+|\Delta|\), with
\(|\Delta|=0,0.001,0.002,0.005,0.01,0.02\) nats. Top: 0--1; bottom:
0--0.15. Each point has 20,000 replicates.}
\label{fig:imbalance-q}
\resultinterpretation{At zero difference, Wald is near or above 0.05 while
Expanded Welch is below it in both allocations. When \(Q\) has more observations
\((n_P,n_Q)=(50,250)\), rejection rises with \(I(Q)-I(P)\); at 0.02 nats it
reaches \(\ImbalanceQHigherWaldRate\) for Wald and
\(\ImbalanceQHigherExpandedRate\) for Expanded Welch. With \(P\) larger, the
curves rise much less. These points are valid, but the alternative gap is not
a size-adjusted power comparison because their matched null rates differ.}
\end{figure}

\begin{figure}[htbp]
\centering
\includegraphics[width=0.96\textwidth]{imbalance_3x3_p_higher.pdf}
\caption[Unequal samples, \(P\) higher]{Unequal samples, \(P\) higher. Columns have
\((n_P,n_Q)=(50,250)\) and \((250,50)\). \(P\)'s \(3\times3\) row and
column margins are \((0.7,0.15,0.15)\); \(Q\)'s are \((0.8,0.1,0.1)\).
Both use the first-cell change block. \(I(Q)=0.02\) and
\(I(P)=0.02+|\Delta|\), with
\(|\Delta|=0,0.001,0.002,0.005,0.01,0.02\) nats.
Top: 0--1; bottom: 0--0.15. Each point has 20,000 replicates.}
\label{fig:imbalance-p}
\resultinterpretation{The zero-difference rates are the same as in
Figure~\ref{fig:imbalance-q}, because each allocation uses the same pair of
populations in both figures at that point; \(P\) and \(Q\) still have different
margins. Reversing the MI direction changes what happens next: when \(P\) has
the larger sample, Wald rises well above its null rate, whereas it stays
nearly flat when \(Q\) has the larger sample. At \((n_P,n_Q)=(50,250)\) and
0.02 nats, Wald rejects
\(\ImbalancePHigherWaldRate\) and Expanded Welch
\(\ImbalancePHigherExpandedRate\); both calculations are valid. Direction and
allocation therefore matter even at the same absolute MI difference.}
\end{figure}

Together, these two figures show why averaging across the directions of
sample imbalance would hide a real difference in performance. Neither
figure alone establishes a method's power advantage at matched false-positive
rates.

\section{Larger differences and rectangular tables}

The main grid deliberately concentrates on small MI differences. A separate
\(2\times2\) family extends the feasible axis to 0.2 nats without changing
the axis units, using only targets the construction can reach. This shows where
detection eventually rises, while retaining the zero-difference point that
defines null calibration. Rectangular \(2\times3\) and \(3\times5\) tables
check that the discussion is not tied only to square alphabets. The panels
do not average over shapes.

\begin{figure}[htbp]
\centering
\includegraphics[width=0.92\textwidth]{broad_2x2_uniform.pdf}
\caption[Wider MI differences, uniform \(2\times2\)]{Wider MI differences in a uniform \(2\times2\) population pair.
The single panel has \(P\) and \(Q\) row and column margins \((0.5,0.5)\),
first-cell changes, \(n_P=n_Q=100\), \(I(P)=0.02\), and
\(I(Q)=0.02+|\Delta|\). The horizontal grid is
\(|\Delta|=0,0.001,0.002,0.005,0.01,0.02,0.05,0.1,0.2\) nats;
the rejection axis is 0--1. Each point has 20,000 replicates.}
\label{fig:broad-focus}
\resultinterpretation{At zero difference, both methods are below the 0.05
line; Expanded Welch is lower, and validity is high. Rejection then rises markedly as the true MI
difference reaches the larger targets: by 0.2 nats both methods reject more
than 90\% of samples. Expanded Welch remains below Wald at intermediate
differences, but the gap narrows when detection is already high. This wider
axis shows that the low rejection over the main grid's small differences is
not a permanent ceiling on detection.}
\end{figure}

\FloatBarrier

\begin{figure}[htbp]
\centering
\includegraphics[width=0.94\textwidth]{rectangular_different_skew.pdf}
\caption[Rectangular tables]{Rectangular tables. Columns are \(2\times3\) and \(3\times5\).
In each dimension \(k\), \(P\)'s corresponding margin is \(m_k(0.7)\)
and \(Q\)'s is \(m_k(0.8)\), for both rows and columns. Thus the first
column uses \(P\) rows \((0.7,0.3)\), columns
\((0.7,0.15,0.15)\), and \(Q\) rows \((0.8,0.2)\), columns
\((0.8,0.1,0.1)\); the second applies the same formula to 3 and 5
categories. Both use first-cell changes, \(n_P=n_Q=100\),
\(I(P)=0.02\), and \(I(Q)=0.02+|\Delta|\) for
\(|\Delta|=0,0.001,0.002,0.005,0.01,0.02\) nats.
Top: 0--1; bottom: 0--0.15. Each point has 20,000 replicates.}
\label{fig:rectangular-focus}
\resultinterpretation{Both rectangular panels start below 0.05, with
Expanded Welch lower. Rejection grows only modestly by 0.02 nats and remains
below 0.05 for both methods throughout the displayed range; validity is high.
Thus these shapes show weak detection over the chosen differences, not an
improvement in false-positive calibration from the lower Expanded Welch
curve. The conservative pattern is not confined to square tables.}
\end{figure}

\FloatBarrier

The wider \(2\times2\) curve shows where detection eventually becomes strong;
the rectangular panels show that the conservative pattern persists in other
table shapes over smaller differences.

\section{Comparing Two Population Constructions}

The additive block is interpretable, but it is only one way to make tables
with fixed margins and target MI\@. Figure~\ref{fig:construction-focus}
compares it with an ordinal log-linear construction under matched targets.

\begin{figure}[htbp]
\centering
\includegraphics[width=0.94\textwidth]{construction_5x5_different_skew.pdf}
\caption[Additive versus log-linear construction]{Additive versus ordinal log-linear construction for
\(5\times5\) tables. The left column uses the first-cell additive block;
the right fits an ordinal interaction while preserving the margins. In both,
\(P\)'s row and column margins are
\((0.7,0.075,0.075,0.075,0.075)\) and \(Q\)'s are
\((0.8,0.05,0.05,0.05,0.05)\). Both have
\(I(P)=0.02\), \(I(Q)=0.02+|\Delta|\), \(n_P=n_Q=100\), and
\(|\Delta|=0,0.001,0.002,0.005,0.01,0.02\) nats.
Top: 0--1; bottom: 0--0.15. Each point has 20,000 replicates.}
\label{fig:construction-focus}
\resultinterpretation{At zero difference, the additive and log-linear Wald
rates are \(\ConstructionAdditiveWaldRate\) and
\(\ConstructionLoglinearWaldRate\); Expanded Welch gives
\(\ConstructionAdditiveExpandedRate\) and
\(\ConstructionLoglinearExpandedRate\). All four are valid and below 0.05.
Rejection stays low or falls slightly as the displayed difference grows.
Matching margins and MI therefore does not make finite-sample behaviour
identical across constructions, but neither construction removes the
conservative gap between the methods.}
\end{figure}

\section{Extreme skew, independence and larger samples}

The extreme-skew check deliberately keeps configurations with very small
expected cell counts. Figure~\ref{fig:extreme-focus} compares three specified
skew levels at \(n_P=n_Q=100\).

\begin{figure}[htbp]
\centering
\includegraphics[width=0.96\textwidth]{extreme_3x3.pdf}
\caption[Extreme-skew \(3\times3\) tables]{Extreme-skew \(3\times3\) tables. Across columns the dominant
row and column probabilities for \((P,Q)\) are
\((0.9,0.95)\), \((0.99,0.995)\), and \((0.999,0.9995)\). Each remaining
margin entry splits the leftover probability equally between the two other
categories. Both use first-cell changes, \(n_P=n_Q=100\),
\(I(P)=0.00001\), and \(I(Q)=I(P)+|\Delta|\), for
\(|\Delta|=0,0.00001,0.00005,0.0001,0.0002\) nats.
Top: 0--1; bottom: 0--0.15. Each point has 20,000 replicates; hollow
markers flag validity below 90\%.}
\label{fig:extreme-focus}
\resultinterpretation{For dominant margins 0.9 and 0.95, the null rates are
0.0701 for Wald and 0.0088 for Expanded Welch: the adjustment moves an
excessive rate past 0.05. At the most extreme margins, 0.999 and 0.9995,
both rates are zero, but only 0.0118 of Wald and 0.00005 of Expanded Welch
null calculations are valid. Rejection does not meaningfully grow over the
displayed alternatives there. The hollow markers show why near-zero
unconditional rejection cannot be read as good calibration in sparse data.}
\end{figure}

Exact independence is a separate theoretical boundary.
Figure~\ref{fig:independence-focus} starts \(P\) at MI zero and includes the
equal-MI point where both populations are independent.

\begin{figure}[htbp]
\centering
\includegraphics[width=0.96\textwidth]{independence_2x2.pdf}
\caption[Exact-independence boundary]{Exact-independence boundary for \(2\times2\) tables. Columns are
uniform margins \((0.5,0.5)\) for both populations, same skew
\((0.8,0.2)\) for both, and different skew with \(P=(0.7,0.3)\) and
\(Q=(0.8,0.2)\); each vector is used for both rows and columns. The
first-cell block is used, \(I(P)=0\), \(I(Q)=|\Delta|\),
\(n_P=n_Q=100\), and
\(|\Delta|=0,0.001,0.002,0.005,0.01,0.02\) nats.
Top: 0--1; bottom: 0--0.15. Each point has 20,000 replicates.}
\label{fig:independence-focus}
\resultinterpretation{At the uniform zero-difference point, Wald rejects
0.00025 and Expanded Welch 0.00010 of samples; validity is high. Rejection grows little over
most of the displayed alternative range, and Expanded Welch remains lower.
This is not a routine positive-MI validation: at exact independence the
population first-order MI variance is zero, so the approximation underlying
both curves changes. Appendix~\ref{app:independence-boundary} explains the
different reference used by a one-table independence test.}
\end{figure}

Does increasing \(n\) resolve the null error when the populations remain
fixed? Figure~\ref{fig:convergence-focus} keeps equal MI and changes only
sample size. It is a null-only comparison, not a power curve.

\begin{figure}[htbp]
\centering
\includegraphics[width=0.9\textwidth]{convergence.pdf}
\caption[Large-sample null behaviour]{Large-sample null behaviour. Rows compare \(2\times2\) and
\(8\times8\) tables; columns have common true MI 0.0001 and 0.02 nats.
For both shapes, \(P\)'s row and column margins are \(m_k(0.7)\) and
\(Q\)'s are \(m_k(0.8)\), with \(k=2\) or 8. Both populations use the
first-cell block. The horizontal axis is the equal sample size
\(n_P=n_Q=1{,}000,2{,}500,10{,}000,50{,}000\) on a labelled log scale;
the vertical axis is null rejection from 0 to 0.10. Each point has 20,000
replicates.}
\label{fig:convergence-focus}
\resultinterpretation{The horizontal axis is sample size, and every point
is a zero-difference null; there are no alternative-detection points. All
displayed calculations are valid. In the
displayed \(2\times2\) and \(8\times8\) settings, rates move toward 0.05
more clearly when common MI is 0.02 than when it is 0.0001. Across all 12
first-block shape--margin combinations at \(n=50{,}000\), the 0.02-MI
null rates range from 0.0470 to 0.0527 for Wald and 0.0466 to 0.0526 for
Expanded Welch. At MI 0.0001, the ranges remain 0.0089--0.0280 and
0.0076--0.0138. Convergence can therefore be slow near independence.}
\end{figure}

\section{Null Rejection Rates and the Standard-Error Check}
\label{sec:mechanism-results}

The main-grid null contains two cases distinguished by whether the population
tables themselves are equal. The 72 uniform and same-skew configurations have
\(P=Q\), whereas the 36 different-skew configurations have \(P\ne Q\) but
\(I(P)=I(Q)\). These cases, where the tables differ but their MI values
agree, most directly match the thesis motivation. Table~\ref{tab:null-band-summary} reports these
profiles separately. The broad interval 0.025--0.075 is used only to describe
the landscape; it is not a formal acceptance region for calibration.

\begin{table}[htbp]
\centering
\small
\caption{Main-grid null configurations below, inside or above the descriptive rejection-rate interval 0.025--0.075.}
\label{tab:null-band-summary}
\begin{tabular}{@{}llrrrr@{}}
\toprule
Margins & Method & Below & Inside & Above & Validity below 90\%\\
\midrule
Uniform & Normal Wald & 10 & 18 & 8 & 4 \\
Uniform & Expanded Welch & 19 & 16 & 1 & 5 \\
Same skew & Normal Wald & 14 & 15 & 7 & 8 \\
Same skew & Expanded Welch & 25 & 11 & 0 & 12 \\
Different skew & Normal Wald & 9 & 16 & 11 & 8 \\
Different skew & Expanded Welch & 19 & 13 & 4 & 12 \\
\bottomrule
\end{tabular}
\end{table}


\FloatBarrier

The table makes the direction of the adjustment clear. Expanded Welch greatly
reduces the number of rates above the descriptive band, but it also increases
the number below it in every margin profile. Its apparent protection against
liberal rejection is therefore purchased with additional conservatism rather
than with calibration in both directions.

A separate supplementary study tested the mechanism without altering the
confirmatory results. For the representative \(3\times3\) different-skew
null with \(I(P)=I(Q)=0.02\) and \(n=100\), ordinary Wald divided each
observed bias-corrected MI difference by its own plug-in standard error
and rejected 0.0166 of sample pairs. Keeping those same differences but dividing
each by one fixed finite-sample standard deviation estimated from an
independent pilot gave a rejection rate of 0.0551. Dividing instead by
the fixed population first-order standard deviation gave 0.0976. The
confirmatory Wald rate was 0.0171; its small difference from 0.0166 is
Monte Carlo variation between independent runs.

\pagebreak[4]
Across the evaluated pairs, the average squared plug-in standard error was
about 1.17 times the finite-sample variance of the MI difference estimated
from the independent pilot. The population
first-order variance was smaller than that finite-sample variance. These
differences help explain why the three denominators produce different
rejection rates. But replacing each sample's standard error with a fixed
number also removes its variation and its dependence on the estimated MI
difference. The contrast between the rejection rates cannot be assigned to
variance inflation alone or split into separate shares.
Section~\ref{sec:mi-variance-covariance}
shows why each population's MI and variance estimates can be related. The
two fixed denominators are diagnostics, not proposed tests: Expanded Welch
changes neither the observed MI differences nor their plug-in standard
errors. Full pilot and method-comparison results are retained in the
experiment supplement.

\section{Runtime and summary of evidence}
\label{sec:runtime-summary}

The full method calls have the same \(O(rc)\) asymptotic complexity, as shown
by the operation count in Section~\ref{sec:expanded-complexity}. Table~\ref{tab:runtime-by-shape}
gives scalar-call times for five shapes; the saved runtime file in
Appendix~\ref{app:evidence-index} contains all 60 timing regimes.

\begin{table}[htbp]
\centering
\small
\caption{Scalar-call times in microseconds for the different-skew
null at \(n_P=n_Q=100\) and \(I(P)=I(Q)=0.02\). Entries give median
[interquartile range] over five passes on 200 saved inputs per regime.
Population construction and multinomial sampling are excluded.}
\label{tab:runtime-by-shape}
\begin{tabular}{@{}lccc@{}}
\toprule
Shape & Normal Wald & Expanded Welch & Expanded/Wald\\
\midrule
\(2\times2\) & 85.8 [85.5, 86.3] & 150.0 [149.4, 151.3] & 1.75\\
\(3\times3\) & 85.1 [84.6, 85.6] & 150.0 [148.9, 150.9] & 1.76\\
\(3\times5\) & 85.2 [84.9, 85.5] & 150.1 [149.5, 150.9] & 1.76\\
\(5\times5\) & 85.3 [85.1, 85.7] & 149.2 [148.5, 150.1] & 1.75\\
\(8\times8\) & 88.0 [87.1, 88.9] & 153.4 [152.4, 154.8] & 1.74\\
\bottomrule
\end{tabular}
\end{table}

\FloatBarrier

Across the 60 regimes, median times are about 85 microseconds for Wald and
150 for Expanded Welch. The median ratio is about 1.75, with regime ratios
about 1.72--1.79. In the displayed \(3\times3\) regime, the medians are
\(\RuntimeRegularWaldMedian\) and \(\RuntimeRegularExpandedMedian\)
microseconds. These machine-specific timings mainly measure fixed per-call
overhead, not the \(O(rc)\) scaling itself. The extra work is feasible but
does not supply a reason to prefer Expanded Welch when the statistical results
show no general improvement.

The evidence gives three conclusions. Expanded Welch can move a liberal Wald
null rate closer to 0.05, but its heavier-tailed reference also reduces
rejection where Wald is already conservative. In sparse samples, the
unconditional rate can fall because Expanded Welch is more often invalid under
the strict component-validity rule. Some failures would not occur under the
direct combined formula and are not needed for the first two findings. These
behaviours cannot be summarised by one number; the matched curves and validity
records underpin the discussion in Chapter~\ref{chap:discussion}.
